% ===== PRÉAMBULE CANONIQUE — à coller VERBATIM en tête de CHAQUE .tex =====
\documentclass[11pt,a4paper]{article}
\usepackage[utf8]{inputenc}\usepackage[T1]{fontenc}\usepackage{lmodern}
\usepackage[french]{babel}
\usepackage{amsmath,amssymb,mathtools}\usepackage{icomma}
\usepackage[version=4]{mhchem}          % \ce{...} : équations & formules
\usepackage{siunitx}\sisetup{locale=FR} % \qty{}{}, \si{}, \ang{}
\usepackage{chemfig}                    % \chemfig{...} : structures développées
\usepackage{xcolor}\usepackage{colortbl}
\usepackage{tikz}\usepackage{pgfplots}\pgfplotsset{compat=1.18}
\usepackage[most]{tcolorbox}\usepackage{enumitem}\usepackage{array,booktabs}
\usepackage[a4paper,margin=2.2cm]{geometry}
\usetikzlibrary{arrows.meta,calc,angles,quotes,positioning,patterns,backgrounds,shapes.geometric}
\definecolor{primary}{RGB}{222,49,99}\definecolor{primarydark}{RGB}{150,22,55}
\definecolor{defcol}{RGB}{0,114,178}\definecolor{methcol}{RGB}{52,92,148}
\definecolor{excol}{RGB}{0,135,90}\definecolor{attncol}{RGB}{200,40,55}
\tcbset{boxbase/.style={enhanced,breakable,boxrule=0.9pt,arc=2.5pt,
  left=3mm,right=3mm,top=2.3mm,bottom=2.3mm,fonttitle=\bfseries,coltitle=white}}
% --- boîtes TUTORIEL (noms EXACTS, ne pas renommer) ---
\newtcolorbox{cle}[1][]{boxbase,colback=defcol!6,colframe=defcol,
  title={\textsc{À retenir}\ifx\relax#1\relax\relax\else\textmd{ : \itshape #1}\fi}}
\newtcolorbox{methode}[1][]{boxbase,colback=methcol!7,colframe=methcol,sharp corners,
  title={\textsc{Méthode}\ifx\relax#1\relax\relax\else\textmd{ --- \itshape #1}\fi}}
\newtcolorbox{exemplebox}[1][]{boxbase,colback=excol!5,colframe=excol,
  title={\textsc{Exemple}\ifx\relax#1\relax\relax\else\textmd{ : \itshape #1}\fi}}
\newtcolorbox{piege}[1][]{boxbase,colback=attncol!6,colframe=attncol,
  title={\textsc{Piège}\ifx\relax#1\relax\relax\else\textmd{ : \itshape #1}\fi}}
\newtcolorbox{remarque}[1][]{boxbase,colback=black!4,colframe=black!45,
  title={\textsc{Remarque}\ifx\relax#1\relax\relax\else\textmd{ : \itshape #1}\fi}}
% --- boîtes EXERCICE / CORRIGÉ (noms EXACTS, auto-comptées) ---
\newtcolorbox[auto counter]{exo}[1][]{boxbase,colback=primary!4,colframe=primary,
  title={\textsc{Exercice}~\thetcbcounter\ifx\relax#1\relax\relax\else\textmd{ --- \itshape #1}\fi}}
\newtcolorbox[auto counter]{corrige}[1][]{boxbase,colback=excol!4,colframe=excol,
  title={\textsc{Corrigé}~\thetcbcounter\ifx\relax#1\relax\relax\else\textmd{ --- \itshape #1}\fi}}
\newcommand{\R}{\mathbb{R}}\newcommand{\N}{\mathbb{N}}\newcommand{\Z}{\mathbb{Z}}\newcommand{\ds}{\displaystyle}
% =========================================================================
\begin{document}
% THEME_TITLE: Influence du pH sur la solubilité
\sisetup{per-mode=symbol}
\section*{\color{primarydark}Thème 5 — Influence du pH sur la solubilité}

Beaucoup de sels peu solubles mettent en jeu des ions dont la concentration dépend du pH (\ce{OH-}, ou
la base conjuguée d'un acide faible). Le pH devient alors un \textbf{levier} pour dissoudre ou faire
précipiter. La synthèse du livre distingue \textbf{trois cas}.

\begin{cle}[{lien entre $[\ce{OH-}]$ et le pH (à \qty{25}{\degreeCelsius})}]
Le produit ionique de l'eau donne $K_e=[\ce{H3O+}][\ce{OH-}]=\num{1.0e-14}$, d'où
\[ \boxed{\,\mathrm{pH} = 14 + \log[\ce{OH-}]\,} \qquad\text{et}\qquad
   \boxed{\,[\ce{OH-}] = 10^{\,\mathrm{pH}-14}\,} \]
Plus le pH est élevé (basique), plus $[\ce{OH-}]$ est grand.
\end{cle}

\begin{cle}[les trois cas à connaître]
\begin{itemize}[leftmargin=1.6em,topsep=2pt,itemsep=2.5pt]
  \item \textbf{Cas 1 --- Hydroxydes \ce{M(OH)_n}} : leur $K_{\mathrm{ps}}$ contient $[\ce{OH-}]^{n}$.
        Quand \textbf{pH $\uparrow$}, $[\ce{OH-}]\uparrow$ et ils \textbf{précipitent} (solubilité
        $\downarrow$). Ils se \emph{dissolvent} en milieu acide.
  \item \textbf{Cas 2 --- Sels d'acides faibles} (\ce{CaCO3}, \ce{FeS}, \ce{CaC2O4}, \ce{CaF2}\dots) :
        l'anion est la base conjuguée d'un acide faible. En milieu acide (\textbf{pH $\downarrow$}),
        \ce{H3O+} consomme l'anion : la solubilité \textbf{augmente}.
  \item \textbf{Cas 3 --- Autres sels} (\ce{AgCl}, \ce{BaSO4}, \ce{PbBr2}\dots) : l'anion vient d'un
        acide \emph{fort}, il ne réagit pas avec \ce{H3O+}. Solubilité \textbf{indépendante} du pH.
  \end{itemize}
\end{cle}

\begin{methode}[pH de début de précipitation d'un hydroxyde]
La précipitation de \ce{M(OH)_n} débute quand $Q=K_{\mathrm{ps}}$, soit
$[\ce{M^{n+}}][\ce{OH-}]^{n}=K_{\mathrm{ps}}$. On isole $[\ce{OH-}]$ puis on passe au pH :
\[ [\ce{OH-}] = \sqrt[n]{\dfrac{K_{\mathrm{ps}}}{[\ce{M^{n+}}]}}
   \quad\Longrightarrow\quad
   \boxed{\,\mathrm{pH}_{\text{précip.}} = 14 + \log\!\sqrt[n]{\dfrac{K_{\mathrm{ps}}}{[\ce{M^{n+}}]}}\,} \]
\end{methode}

\begin{exemplebox}[début de précipitation de \ce{Al(OH)3}]
$[\ce{Al^{3+}}]=\qty{1.0e-3}{\mole\per\litre}$, $K_{\mathrm{ps}}=\num{2.0e-32}$ :
\[ [\ce{OH-}]=\sqrt[3]{\dfrac{\num{2.0e-32}}{\num{1.0e-3}}}=\sqrt[3]{\num{2.0e-29}}\approx\qty{2.7e-10}{\mole\per\litre}, \]
\[ \mathrm{pH}=14+\log(\num{2.7e-10})\approx 4{,}4. \]
\ce{Al(OH)3} précipite dès que l'eau cesse d'être franchement acide (floculation dans le traitement des
eaux).
\end{exemplebox}

\begin{exemplebox}[le calcaire se dissout dans l'acide (cas 2)]
En milieu acide, \ce{CaCO3} se dissout par une réaction \emph{totale} :
\[ \ce{CaCO3(s) + 2 H3O+ -> Ca^{2+}(aq) + CO2(g) + 3 H2O}. \]
L'ion \ce{CO3^{2-}} est consommé et transformé en \ce{CO2} gazeux : une goutte de vinaigre sur de la
craie fait effervescence, l'eau pure non.
\end{exemplebox}

\begin{piege}[pH ou pOH ?]
Pour un hydroxyde, ne pas confondre : $\mathrm{pH}=14+\log[\ce{OH-}]$ (et non $-\log[\ce{OH-}]$, qui
serait le pOH). Rappel : $\mathrm{pH}+\mathrm{pOH}=14$.
\end{piege}

\begin{remarque}[direction à retenir]
Hydroxydes : pH $\uparrow\Rightarrow s\downarrow$. Sels d'acides faibles : pH $\downarrow\Rightarrow
s\uparrow$. Autres : aucun effet. Dans tous les cas \emph{sensibles}, le sel est \textbf{plus soluble en
milieu acide}.
\end{remarque}
\end{document}
